%% kapa: exam=1 round=23 year=2012 period=1 form=A subject=civil_law subject_ko=민법 seq=35 number=35
%% source: Q-Net 감정평가사 기출문제 (cst00309) artlSeq=5052578 file=raw/1008/5052578/2012년 제23회 감정평가사 제1차시험 기출문제(1교시 A형).pdf
%% conversion: kordoc 4.24.51 -> tools/gen1.py (auto); review=proofread-sonnet; flags=none
\begin{question}{35}
저당권의 침해에 관한 설명으로 옳지 않은 것은? (다툼이 있으면 판례에 의함)
\begin{choices}
  \item 저당권이 설정된 임야의 수목이 부당하게 벌채ㆍ반출된 경우, 저당권자는 자신에게 수목의 반환을 청구할 수 있다.
  \item 선순위저당권이 소멸된 경우, 후순위저당권자는 선순위저당권의 말소등기를 청구할 수 있다.
  \item 제3자가 저당권의 목적물을 손상시켜도 나머지 부분의 가액이 피담보채권을 담보하기에 충분한 경우에는 저당권자는 그 제3자에 대하여 손해배상을 청구할 수 없다.
  \item 저당권설정자의 책임 있는 사유로 인하여 저당물의 가액이 현저히 감소된 때에는 저당권자는 원상회복 또는 상당한 담보제공을 청구할 수 있다.
  \item 특별한 사정이 없는 한 저당 토지의 지상권자가 통상의 용법에 따라 토지를 사용ㆍ수익하는 경우, 저당권 침해에 해당하지 않는다.
\end{choices}
\end{question}
\answer{1}
